Training jobs on preemptible machines get a notice of thirty seconds to two minutes before the machine is reclaimed.
Every published elastic or spot training system uses that window to move the complete optimizer state, which for Adam is
three times the size of the model, and treats the cost of recovery as a constant. This paper measures which parts of
that state have to cross. A training run is interrupted at a chosen step, resumed from \nArms{} different payloads, and
compared with a bit-identical uninterrupted control over \nCells{} migrations, three model sizes and up to \nSeedsMax{}
seeds, against a noise floor set by the spread between seeds. Three results follow. Quantised moments at half the bytes
are indistinguishable from the full state. A row-and-column sketch of the second moment, with the first moment dropped,
cuts the payload to a third, or a sixth with half-precision weights, and stays inside the noise floor. Dropping both
moments does not, and the damage is concentrated early in training; the two repairs in common use, resetting the step
counter and re-estimating the moments from fresh gradients, are the two worst options measured. A short analysis
predicts the size and duration of each transient from the two decay rates alone and agrees with measured update sizes.
Payload sizes are measured from the serialised bytes, drain and restore times are timed stage by stage, and the same
payloads are resumed on a CPU, a discrete GPU and an integrated GPU. A cost model turns the measurements into a rule
for what to send inside a given notice window.
